\chapter{Methods}

\section{Overview of the Pipeline}
\label{sec:methods_overview}
This chapter describes the end-to-end methodology used to build an explainable binary classifier for anticipating \textit{rapid} versus \textit{slow} functional progression in amyotrophic lateral sclerosis (ALS) patients, using PRO-ACT data.
The pipeline follows a strict leakage-free design:
(i) define a subject-level baseline timepoint $t_0$,
(ii) construct predictors using only information available up to $t_0$,
(iii) define targets using outcomes observed after $t_0$ for fixed horizons,
and (iv) evaluate models using subject-wise splits to avoid overlap between training and test subjects.
Figure~\ref{fig:pipeline_placeholder} can be used to summarize the stages: baseline definition, target construction, feature extraction, model training/evaluation, ablation analysis, hyperparameter optimization, and explainability (XAI).

\begin{figure}[ht]
  \centering
  \fbox{\parbox{0.9\linewidth}{\centering Placeholder: pipeline overview diagram}}
  \caption{End-to-end pipeline (baseline $\rightarrow$ targets $\rightarrow$ features $\rightarrow$ training/evaluation $\rightarrow$ ablation $\rightarrow$ Optuna tuning $\rightarrow$ threshold sanity check $\rightarrow$ SHAP/XAI).}
  \label{fig:pipeline_placeholder}
\end{figure}

\section{Dataset}
\label{sec:methods_dataset}
\subsection{PRO-ACT Data Sources}
\label{sec:methods_proact_sources}
We use the Pooled Resource Open-Access ALS Clinical Trials (PRO-ACT) database, the largest publicly available repository of ALS clinical trial data, aggregating records from multiple completed trials.
The raw database comprises 16 tables with longitudinal clinical, demographic, and laboratory measurements.
In this work, we employ the following source tables:
\begin{itemize}
  \item \textbf{ALSFRS-R:} longitudinal functional assessments (used for baseline anchoring and target construction),
  \item \textbf{Demographics:} age, sex, ethnicity, and race,
  \item \textbf{Vital Signs:} weight, height, BMI, pulse, respiratory rate, temperature, and blood pressure variants (systolic, diastolic, standing, supine),
  \item \textbf{FVC:} forced vital capacity measurements (litres, percent of normal) as a marker of respiratory function,
  \item \textbf{Riluzole:} administration records of riluzole (the only disease-modifying therapy approved during the trial period),
  \item \textbf{Concomitant Medications:} count of concurrent medications as a proxy for comorbidity burden,
  \item \textbf{Labs:} serum creatinine and alanine aminotransferase (ALT), selected due to their established association with ALS prognosis in prior literature~\cite{Anani2023}.
\end{itemize}
We deliberately restrict the feature set to routinely collected clinical variables to maximise translational potential; imaging, genetic, and electrophysiological markers are excluded because they are not universally available in clinical practice.

\subsection{Study Population and Eligibility}
\label{sec:methods_population}
From the full PRO-ACT database ($N\!=\!5{,}436$ subjects with at least one ALSFRS-R record), subjects are retained if they satisfy three conditions:
(i) a valid ALSFRS-R total score at baseline ($t_0$),
(ii) at least one further ALSFRS-R measurement within the target-horizon window (Section~\ref{sec:methods_targets}), and
(iii) a computable functional slope that meets the quality threshold (Section~\ref{sec:methods_slope}).
These criteria yield horizon-specific cohorts: $N\!=\!2{,}407$ for 3 months and $N\!=\!1{,}392$ for 6 months, with 720 subjects eligible at both horizons.
The reduction from 5{,}436 to 1{,}392 (6-month) is primarily driven by the requirement for follow-up measurements near the 180-day mark.

\subsection{Held-Out Data Partitioning}
\label{sec:methods_holdout}
Before any modelling, each horizon-specific cohort is partitioned into a \textbf{development (DEV)} set (80\%) and a \textbf{held-out test (TEST)} set (20\%) using a stratified random split (stratified by the binary target to preserve class proportions).
For the 6-month horizon this yields DEV $n\!=\!1{,}113$ and TEST $n\!=\!279$; for 3 months, DEV $n\!=\!1{,}925$ and TEST $n\!=\!482$.
All hyperparameter tuning, cross-validation, ablation, and model selection are performed exclusively on the DEV set.
The TEST set is reserved for a single, unbiased evaluation of the final selected model (Section~\ref{sec:results_test_eval}), preventing any form of selection bias or inadvertent data snooping.

\section{Baseline Definition and Leakage-Free Protocol}
\label{sec:methods_baseline}
\subsection{Baseline Timepoint ($t_0$)}
\label{sec:methods_t0}
For each subject, the baseline timepoint $t_0$ is defined as the \textbf{first visit with a valid ALSFRS-R total score}.
This choice provides a consistent subject-level anchor and maximizes cohort size while remaining clinically meaningful.
Using the first visit rather than a later one avoids selecting subjects with atypically long follow-up, which could bias the cohort towards slower progressors.

\subsection{Why Baseline Matters}
\label{sec:methods_baseline_why}
The baseline timepoint is critical to ensure temporal validity and prevent information leakage.
All predictors (features) must be constructed using data available \emph{up to and including} $t_0$, while all targets must be derived exclusively from measurements \emph{after} $t_0$.
This strict temporal separation mirrors the intended clinical use case: a clinician wishes to assess, at a patient's first structured evaluation, whether the patient is likely to experience rapid functional decline over the following months.
Violating this separation---for example, by including post-baseline lab values as predictors---would produce optimistically biased models that cannot be replicated at the point of clinical decision-making.

\section{Target Construction: ALSFRS-R Slopes}
\label{sec:methods_targets}
\subsection{Horizon Definition and Temporal Tolerance}
\label{sec:methods_horizons}
We define two prediction horizons:
\begin{itemize}
  \item \textbf{3 months:} $H=90$ days with tolerance $\pm 7$ days,
  \item \textbf{6 months:} $H=180$ days with tolerance $\pm 7$ days.
\end{itemize}
Tolerance accounts for real-world visit variability while preserving temporal specificity.

\subsection{Slope Estimation}
\label{sec:methods_slope}
Functional progression is quantified via the ALSFRS-R slope from baseline to the horizon.
For each subject and horizon, we collect all ALSFRS-R measurements within the window $[t_0,\; t_0 + H \pm 7\text{d}]$ and estimate the rate of functional change by ordinary least squares regression of the ALSFRS-R total score against time (in days).
The resulting coefficient is normalised to \emph{points per 30\,days} to provide a clinically intuitive unit.
A negative slope indicates functional decline; the more negative the slope, the faster the patient is deteriorating.

Slope quality is monitored via two indicators:
(i) the number of ALSFRS-R measurements contributing to each estimate ($n_{\text{points}}$), and
(ii) the effective time span covered.
For the 6-month cohort, 99.8\% of slope estimates are based on $\geq 3$ data points (median $n_{\text{points}} = 7$), indicating robust estimation.
For 3 months, 90.5\% of slopes use $\geq 3$ points (median $n_{\text{points}} = 4$), reflecting the shorter observation window.

\subsection{Binary Target: Rapid vs Slow Progression}
\label{sec:methods_binary_target}
The classification target is defined as \textit{rapid} versus \textit{slow} progression based on slope severity.
We define the positive class (\textit{rapid}) as the worst 30\% of slopes (most negative, i.e.\ fastest decline) \textbf{within the training data of each cross-validation fold} (see Section~\ref{sec:methods_validation}).

The choice of the 30th percentile is motivated by three considerations:
\begin{enumerate}
  \item \textbf{Clinical prevalence:} prior literature reports that approximately one-third of ALS patients exhibit accelerated functional decline within the first year~\cite{Pancotti2022}, making 30\% a clinically grounded cutoff.
  \item \textbf{Leakage prevention:} computing the percentile on the training fold only ensures that no information from held-out subjects influences the label assignment, maintaining strict temporal and subject-level separation.
  \item \textbf{Class balance:} a 30/70 split produces a moderate imbalance that is challenging enough to test model robustness while remaining tractable with standard balancing techniques (Section~\ref{sec:methods_imbalance}).
\end{enumerate}
The fold-wise label definition means that the exact slope cutoff varies slightly across folds, but the semantic interpretation---``clinically concerning rate of decline''---remains consistent.

\section{Feature Engineering at Baseline ($t_0$)}
\label{sec:methods_features}
\subsection{Feature Blocks}
\label{sec:methods_feature_blocks}
Features are extracted using only measurements available up to $t_0$.
We group features into clinically coherent blocks:
\begin{itemize}
  \item \textbf{Baseline block} ($n\!=\!15$): demographic variables (age, sex, ethnicity, race indicators), ALSFRS-R total at baseline, and the three respiratory ALSFRS-R sub-items.
  \item \textbf{Vitals block} ($n\!=\!12$): last pre-baseline vital sign measurements (weight, height, BMI, pulse, respiratory rate, temperature, blood pressure variants).
  \item \textbf{Respiratory block} ($n\!=\!3$): FVC at or before baseline (best-trial litres, percent normal, and subject normal flag).
  \item \textbf{Treatment block} ($n\!=\!3$): riluzole use before $t_0$, clinical-trial study arm, and number of concomitant medications.
  \item \textbf{Labs block} ($n\!=\!2$): creatinine and ALT at baseline.
\end{itemize}
The full feature set (v2) comprises 35 predictors after excluding near-constant columns.
This block structure supports interpretability and enables ablation studies (Section~\ref{sec:methods_ablation}).

\subsection{Temporal Selection Rule (Last Pre-Baseline)}
\label{sec:methods_last_prebaseline}
For longitudinal tables (vitals, FVC), we apply a consistent temporal rule:
for each subject, we select the \textbf{most recent measurement with timestamp $\leq t_0$}.
This ``last pre-baseline'' strategy provides clinically plausible baseline status while avoiding leakage.

\subsection{Coverage and Missingness}
\label{sec:methods_coverage}
We quantify feature availability per block and report coverage as the fraction of baseline subjects with non-missing values for each feature.
At the block level, vital signs are available for 81.6\% of subjects, while FVC is available for 53.3\%.
Within the baseline block, ALSFRS-R total and sex have 100\% coverage; age is available for 87.1\% of subjects.
The treatment block has near-complete coverage (riluzole and concomitant medication count are available for all subjects; study arm for 91.1\%).
Labs (creatinine, ALT) are available for 64.7\% and 69.0\% of subjects, respectively.

This heterogeneous coverage motivates the use of median imputation within the preprocessing pipeline (Section~\ref{sec:methods_preprocessing}) and the MAR analysis (Section~\ref{sec:results_fvc_missingness}) to verify that missingness is not outcome-dependent.

\subsection{Data Quality: BMI Correction}
\label{sec:methods_bmi_correction}
During exploratory analysis, we identified 147 raw height records in the vital signs table labelled as ``Centimeters'' that contained values below 120\,cm, consistent with measurements actually recorded in inches but with an incorrect unit label.
These erroneous records produced physiologically implausible BMI values (up to 4{,}512).
To address this, we applied plausible adult bounds: any height outside $[120, 220]$\,cm or weight outside $[30, 200]$\,kg was set to \texttt{NaN}.
This correction affected 132 height and 31 weight records.
After correction, BMI ranges from 14.4 to 46.9 across the cohort.
All v2 datasets used in subsequent analyses were regenerated with the corrected vitals.

\section{Modeling Setup}
\label{sec:methods_modeling}
\subsection{Models Considered}
\label{sec:methods_models}
We evaluate a diverse set of seven supervised learning algorithms spanning different model families, to identify which approaches are best suited for this clinical prediction task:
\begin{itemize}
  \item \textbf{K-Nearest Neighbours (KNN):} a non-parametric instance-based learner; included as a baseline that makes no assumptions about the decision boundary.
  \item \textbf{Decision Tree (DT):} a single tree classifier; included for its intrinsic interpretability and as the building block of ensemble methods.
  \item \textbf{Random Forest (RF):} bagging ensemble of decision trees; reduces variance relative to a single tree and handles mixed feature types well.
  \item \textbf{Support Vector Machine (SVM):} maximum-margin classifier with RBF kernel; effective in moderate-dimensional spaces.
  \item \textbf{Logistic Regression (LR):} linear model; retained as an interpretable baseline whose coefficients provide direct insight into feature effects.
  \item \textbf{LightGBM (LGBM):} gradient-boosted decision trees with histogram-based splits; efficient on medium-sized tabular datasets.
  \item \textbf{XGBoost (XGB):} gradient-boosted trees with regularisation; widely adopted in clinical ML and compatible with SHAP-based post-hoc explainability.
\end{itemize}
This selection covers instance-based, linear, kernel-based, bagging, and boosting paradigms, allowing a principled comparison of model families rather than relying on a single architecture.
Following hyperparameter optimisation (Section~\ref{subsec:methods_optuna}), the final explainable model is selected based on the best cross-validated PR-AUC, with interpretability provided via SHAP (Section~\ref{sec:methods_xai_plan}) and, for Logistic Regression, via direct coefficient analysis.

\subsection{Preprocessing Pipeline}
\label{sec:methods_preprocessing}
To ensure reproducibility and consistent handling of missing and categorical data, we use a unified preprocessing pipeline:
\begin{itemize}
  \item Numeric variables: median imputation; standardization for models sensitive to scaling (LR and SVM),
  \item Categorical variables: most-frequent imputation; one-hot encoding with unseen-category handling.
\end{itemize}
Tree-based models (Decision Tree, Random Forest, LightGBM, XGBoost) do not require scaling, but use consistent imputation and encoding.

\section{Handling Class Imbalance}
\label{sec:methods_imbalance}
The positive class (\textit{rapid}) constitutes approximately 30\% of the cohort, creating a moderate but non-trivial class imbalance.
Although the imbalance ratio ($\approx$1:2.3) is not extreme, it can still bias classifiers towards the majority class, particularly when using accuracy or threshold-dependent metrics at a fixed cutoff of 0.5.
We address this through three complementary strategies:

\begin{enumerate}
  \item \textbf{Metric choice.} Precision--Recall Area Under the Curve (PR-AUC, equivalently Average Precision) is used as the primary ranking and optimisation metric.
    Unlike ROC-AUC, PR-AUC is sensitive to the positive-class prevalence and penalises models that achieve high recall only at the expense of many false positives---a desirable property in a clinical screening context.
    ROC-AUC is additionally reported for comparability with prior ALS literature.
  \item \textbf{Operating-point selection.} Decision thresholds are not fixed at $t\!=\!0.5$.
    Instead, within each cross-validation fold, the threshold that maximises the $F_2$ score on the training data is selected and applied to the fold's test partition.
    $F_2$ weights recall twice as heavily as precision, reflecting the clinical priority of identifying patients at risk of rapid decline (where missing a high-risk patient is worse than a false alarm).
  \item \textbf{Systematic balanced vs.\ unbalanced comparison.} Every model family (except KNN, which has no native class-weighting support) is tuned in two modes:
    \begin{itemize}
      \item \textit{Unbalanced}: the classifier treats both classes equally, with no cost adjustment.
      \item \textit{Balanced}: class weights inversely proportional to class frequencies are applied.
        Concretely, Decision Tree, Random Forest, SVM, and Logistic Regression use \texttt{class\_weight=balanced};
        LightGBM uses \texttt{is\_unbalance=True};
        and XGBoost uses per-fold \texttt{scale\_pos\_weight} computed as $n_{\text{negative}} / n_{\text{positive}}$ on each training fold.
    \end{itemize}
    This yields 13 model$\times$mode configurations (7 models $\times$ 2 modes minus KNN balanced), enabling a controlled comparison of the effect of class weighting on both threshold-free and threshold-dependent metrics.
\end{enumerate}

We deliberately avoid synthetic oversampling techniques (e.g.\ SMOTE, ADASYN) for two reasons:
(i) generating artificial samples in a clinical dataset risks introducing physiologically implausible feature combinations, and
(ii) cost-sensitive learning via class weights achieves a similar effect without altering the training distribution.

\section{Validation Strategy}
\label{sec:methods_validation}
\subsection{Subject-wise Cross-Validation (GroupKFold)}
\label{sec:methods_groupkfold}
All model development (tuning, ablation, model comparison) uses \textbf{GroupKFold} cross-validation with $K\!=\!5$ folds and groups defined by \texttt{subject\_id}.
This guarantees that all records from a given subject appear in exactly one fold---either in training or in test, never both---preventing optimistic bias from subject-level correlations.

Within each fold, the binary target (\textit{rapid} vs \textit{slow}) is re-computed from the training partition's slope distribution (30th percentile), ensuring that label thresholds are never influenced by held-out data.
Performance is reported as mean $\pm$ standard deviation across the five folds, providing a more stable and unbiased estimate than a single fixed split.

\section{Decision Threshold Selection and Sanity Check}
\label{sec:methods_threshold}
\subsection{Fold-wise Threshold Optimization (Maximizing $F_2$)}
\label{sec:methods_threshold_foldwise}
For threshold-dependent metrics (e.g., $F_1$, $F_2$, FP/FN), we do not use a fixed threshold of 0.5.
Instead, within each fold we select a decision threshold on the training data that maximizes $F_2$ and then apply it to the fold test set.
This yields an operating point that prioritizes recall of \textit{rapid} subjects, consistent with a screening-oriented clinical scenario.

\subsection{Final Threshold for the Deployed Model}
\label{sec:methods_threshold_final}
After selecting the final model, we perform a dedicated sanity check to choose a single operating threshold for the final artifact.
We evaluate a grid of thresholds and select a value that preserves high recall while avoiding an unrealistically high predicted-positive rate.
For the tuned XGBoost final model, the selected decision threshold was fixed to $0.25$ for subsequent analyses.

\section{Ablation Study}
\label{sec:methods_ablation}
To quantify the contribution of each feature block, we conduct a cumulative ablation study across six configurations:
\begin{itemize}
  \item A: Baseline only ($n\!=\!15$),
  \item B: A + Vitals ($n\!=\!27$),
  \item C: A + FVC ($n\!=\!18$),
  \item D: A + Vitals + FVC ($n\!=\!30$),
  \item E: D + Treatment ($n\!=\!33$),
  \item F: D + Treatment + Labs ($n\!=\!35$, full v2).
\end{itemize}
Using the \textbf{Optuna-tuned hyperparameters} from the full-feature configuration, each block is evaluated under the same leakage-free CV protocol.
This isolates where predictive signal originates without the confound of re-tuning per block.

\section{Hyperparameter Optimization (Optuna)}
\label{subsec:methods_optuna}
Hyperparameter optimisation is performed using \textit{Optuna}~\cite{optuna2019}, a Bayesian framework based on the Tree-structured Parzen Estimator (TPE) algorithm.
TPE models the objective function as a ratio of two density estimators (one for good trials, one for bad trials) and samples new hyperparameter candidates from the region of parameter space that maximises the expected improvement.
Compared to grid or random search, TPE converges to strong configurations with fewer evaluations, which is important given the computational cost of nested cross-validation.

The search is configured as follows:
\begin{itemize}
  \item \textbf{Objective:} mean PR-AUC across 5 GroupKFold folds (with fold-wise binary labels).
  \item \textbf{Budget:} $N_{\text{trials}}\!=\!60$ per model $\times$ mode configuration.
    With $K\!=\!5$ folds each, this amounts to 300 model fits per configuration and $13 \times 300 = 3{,}900$ fits in total.
  \item \textbf{Reproducibility:} TPE sampler seeded at 42.
  \item \textbf{Hyperparameter spaces:} model-specific ranges covering learning rate, regularisation strength, tree depth, number of estimators, and kernel parameters, as applicable.
\end{itemize}
Additional metrics (ROC-AUC, $F_1$, $F_2$ at threshold 0.5) are recorded for every trial but do not influence the search direction.
The best hyperparameters per configuration are persisted as JSON files for downstream use in ablation and final model training.

This unified protocol---same data splits, same fold-wise label procedure, same number of trials---ensures that all 13 model $\times$ mode configurations are evaluated on an equal footing, eliminating confounds from inconsistent search budgets or differing validation protocols.

\section{Final Model Training and Persistence}
\label{sec:methods_final_model}
Following model selection, we train a final tuned XGBoost pipeline on the full 6-month dataset and persist:
(i) the trained pipeline (preprocessing + classifier) and
(ii) metadata including the feature list, slope cutoff for label definition, and the final decision threshold.
This produces a single reproducible artifact for downstream explainability analyses (Chapter~\ref{chap:results}).

\section{Explainability Plan (XAI)}
\label{sec:methods_xai_plan}
Explainability is treated as a core requirement.
For the final tuned XGBoost model, we use SHAP to provide:
\begin{itemize}
  \item \textbf{Global explanations:} feature importance based on mean absolute SHAP values and summary (beeswarm) plots to characterize the overall contribution of each feature,
  \item \textbf{Local explanations:} per-subject SHAP decompositions (e.g., waterfall plots) for representative TP/FP/FN/TN cases, enabling clinically actionable interpretation of individual predictions.
\end{itemize}
Logistic Regression is retained as an interpretable baseline; its coefficients may be analyzed as a complementary, model-intrinsic explanation of directionality in feature effects.

